\onecolumn
\begin{multicols}{2}
\section{Can CONV Be Made Efficient Using Convolution?}
\label{sec:convstar}

Because CONV specifies the two-jet, lineage ledger, signed-fibre projection,
Bernstein synthesis, and Eikonal factor blend independently of an execution
backend, the same operator admits a bounded-work FIR realization. Denoted by
CONV$^{*}$, this realization eliminates the one-dimensional jet algebra into
fixed convolution banks, batches the exact projection, compiles repeated
phases into a polyphase output bank, and replaces the all-source action sum by
the local nodal compilation derived below. The asterisk therefore denotes a
computational refinement of the same admitted fibre, rather than another
interpolation model. In two dimensions, the FIR bank replaces each
one-dimensional occurrence of \(\mathcal C\) in \(\Tcal_{xy}\) and
\(\Tcal_{yx}\); at a fixed endpoint-aligned scale, the same polyphase bank is
used in opposite factor orders and the resulting proposals are blended
pointwise, without forming a source-specific frame or chart bank.

\subsection{Local compilation of Eikonal admission}

Let \(\eta_p\) be the nodal order coordinate in
\eqref{eq:witnesscompatibility}. A local compilation is obtained by replacing
the all-source action interpolant \(\beta\) with its bilinear nodal extension,
\begin{align}
 \beta^*(x)&=\sum_{p\in\operatorname{cell}(x)}b_p(x)\eta_p,\nonumber\\
 \Tcal^*Y&=(1-\beta^*)\Tcal_{xy}Y+\beta^*\Tcal_{yx}Y.
 \label{eq:starlocaladmission}
\end{align}
\begin{proposition}[Cardinal compilation identity]
Let \(I_h\) denote cellwise bilinear \(Q_1\) interpolation on the source
lattice. Then
\begin{equation}
 \boxed{\beta^*=I_h\beta.}
 \label{eq:starcardinalcompilation}
\end{equation}
Consequently, \(\beta^*\) is the unique continuous cellwise bilinear extension
of the formal Eikonal admission's nodal trace.
\end{proposition}

\begin{proof}
The Kronecker action identity \eqref{eq:kroneckeraction} gives
\(\beta(x_p)=\eta_p\) at every source node. Applying \(I_h\) therefore yields
\((I_h\beta)(x)=\sum_{p\in\operatorname{cell}(x)}b_p(x)\beta(x_p)\), which is
exactly \eqref{eq:starlocaladmission}.
\end{proof}

Thus the compilation is not an independently chosen approximation rule, but
the canonical \(Q_1\) interpolant of the formal admission data; it is
parameter-free, bounded in \([0,1]\), and permutation covariant. Subtracting
\eqref{eq:witnessedblend} from \eqref{eq:starlocaladmission} gives the exact
implementation defect
\begin{equation}
 \Tcal^*Y-\Tcal Y
 = (\beta^*-\beta)(\Tcal_{yx}Y-\Tcal_{xy}Y),
 \label{eq:stardefect}
\end{equation}
and therefore, componentwise,
\begin{equation}
 |\Tcal^*Y-\Tcal Y|
 \leq |\beta^*-\beta|\,|\Tcal_{yx}Y-\Tcal_{xy}Y|.
 \label{eq:stardefectbound}
\end{equation}
Therefore the local compilation is exact wherever the factor orders commute,
including constant and affine coordinate fields. Moreover,
\eqref{eq:stardefect} isolates the sole two-dimensional approximation
introduced by CONV$^{*}$, since the current construction within each factor
remains unchanged.

For direct evaluation, \(\eta_p\) requires only the squared first-jet
channels already produced by the two Cartesian front ends:
\begin{equation}
 \eta_p=
 \frac{\sum_c(D_yY_p^{(c)})^2}
 {\sum_c\bigl[(D_xY_p^{(c)})^2+(D_yY_p^{(c)})^2\bigr]},
 \label{eq:stardirecteta}
\end{equation}
with the value one half when the denominator vanishes. Consequently, neither
an eigendecomposition nor the cross term of the current tensor is required
to evaluate the order coordinate. At a target point, the four nodal values
in \eqref{eq:starlocaladmission} may be combined with the final convex blend
in the same output write. This evaluates \(I_h\eta\) exactly while avoiding a
separate target-sized coordinate raster and its two Cartesian interpolation
passes.

\subsection{Aggregate Green admission}

A global elliptic continuation of the same nodal coordinate is obtained as
follows. Put \(K=M^{-1}\), let \(\rho\) be the source density per unit
Riemannian area, and define
\begin{equation}
 \begin{aligned}
  (\rho-\Delta_M)U&=\sum_a\delta_{s_a},\\
  (\rho-\Delta_M)V&=\sum_a\eta_a\delta_{s_a},\\
  \beta_G&=V/U.
 \end{aligned}
 \label{eq:greenadmission}
\end{equation}
subject to the no-flux boundary condition
\(n^TK\nabla U=n^TK\nabla V=0\). In the unit source-lattice coordinates used
throughout the paper, \(\rho=1\). Under \(x\mapsto cx\), both \(\rho\) and
\(\Delta_M\) scale by \(c^{-2}\), so the screening length is fixed by the
sampling density rather than by a fitted bandwidth. Moreover, the two
right-hand sides use one SPD operator and one factorization, without forming
an individual distance field for each source.

\begin{proposition}[Green transport coordinate]
The continuum coordinate in \eqref{eq:greenadmission} is cardinal by
continuous extension, satisfies
\begin{equation}
 \min_a\eta_a\leq\beta_G(x)\leq\max_a\eta_a,
 \label{eq:greenmaximum}
\end{equation}
and, away from the source sites, obeys
\begin{equation}
 \operatorname{div}\!\left(U^2M^{-1}\nabla\beta_G\right)=0.
 \label{eq:greenharmonic}
\end{equation}
Equivalently, it minimizes
\begin{equation}
 \mathcal E_U[\vartheta]
 =\frac12\int_\Omega U^2\nabla\vartheta^TM^{-1}\nabla\vartheta\,
 d\operatorname{vol}_M
 \label{eq:greenenergy}
\end{equation}
subject to the nodal traces \(\vartheta(s_a)=\eta_a\).
\end{proposition}

\begin{proof}
For \(\rho>0\), the shifted Neumann form is coercive and its Green kernel is
strictly positive. Consequently, \(V/U\) is a positive normalized
Green-kernel blend of the \(\eta_a\), which proves
\eqref{eq:greenmaximum} \cite{gilbargtrudinger2001}. In two dimensions, the
diagonal singularity is
\begin{equation}
 G_\rho(x,s_a)=-\frac{1}{2\pi}\log d_M(x,s_a)+O(1).
 \label{eq:greensingu}
\end{equation}
and therefore the self term dominates at \(s_a\), giving
\(\lim_{x\to s_a}\beta_G(x)=\eta_a\). Away from the sources, substituting
\(V=U\beta_G\) into the two homogeneous resolvent equations cancels the
common mass term; the product rule then yields \eqref{eq:greenharmonic}, whose
weak form is the Euler equation of \eqref{eq:greenenergy}.
\end{proof}

This construction continues the same nodal order coordinate, but it does not
fuse the inverse-square Eikonal admission exactly. Indeed,
\eqref{eq:greensingu} is logarithmic, whereas
\eqref{eq:actionpartition} uses \(d_M^{-2}\); hence no second-order local
elliptic Green function in two dimensions can make the two kernels
identical. Since a finite-stencil diagonal also removes the continuum
singularity, the raw discrete ratio \(V/U\) is not exactly cardinal. Exact
discrete cardinality is obtained by solving \eqref{eq:greenharmonic} with the
nodal traces imposed.

Both the global Green admission and the local compilation therefore have the
trace \(\eta_p\). The latter is the unique cellwise \(Q_1\) continuation
satisfying locality, linearity in the four nodal values, cardinality, and
separate affine exactness, whereas the former is the global minimizer of
\eqref{eq:greenenergy}. Their difference in the reconstructed field is
isolated exactly by
\begin{equation}
 \Tcal_{\beta_G}Y-\Tcal_{\beta^*}Y
 =(\beta_G-\beta^*)(\Tcal_{yx}Y-\Tcal_{xy}Y).
 \label{eq:greenblenddefect}
\end{equation}
Consequently, Green admission supplies a global metric-harmonic coordinate
when that property is required, while \(Q_1\) remains the compact
CONV$^{*}$ realization.

One Cartesian factor has the dataflow
\begin{equation}
 y\xrightarrow{\rm FIR}(a,\delta)
 \xrightarrow{\rm scan}\sigma
 \xrightarrow{\rm tensorized\ projection}c
 \xrightarrow{\rm polyphase}\widehat y.
 \label{eq:starflow}
\end{equation}
Since the factors retain CONV's horizontal--vertical barrier, the second
factor analyzes the values synthesized by the first. No single separable
two-dimensional convolution therefore represents the complete nonlinear
composition.

\subsection{Compact FIR current bank}

On an interior cell, let \(v_i=(y_{i+r})_{r=-2}^{3}\). Eliminating the
intermediate jets and Bezier controls yields
\begin{equation}
 a_i=Kv_i,\quad K=\frac1{240}
 \left[\begin{smallmatrix}
 4&-32&0&32&-4&0\\
 3&-16&-30&48&-5&0\\
 -7&39&-130&130&-39&7\\
 0&5&-48&30&16&-3\\
 0&4&-32&0&32&-4
 \end{smallmatrix}\right].
 \label{eq:starfir}
\end{equation}
Hence the complete interior proposal is a fixed bank of five six-tap FIR
filters. An equivalent realization first computes the two five-tap jet
channels and the secant, and then forms the five currents by pointwise
algebra. The latter requires fewer arithmetic operations on a conventional
CPU, whereas the fused form maps directly to a multi-output convolution
primitive. At the outer and near-outer cells, the stated second-order closure
is retained rather than inferred from \eqref{eq:starfir} or supplied by zero
padding.

\subsection{The irreducible sign scan}

Starting from the sampled signs
\(s_i=\operatorname{sign}(\delta_i)\), CONV$^{*}$ fills a zero from the last
witnessed nonzero sign, using the first future witness only for an initial
zero run. Applying the same local least-disagreement rule to the resulting
transition sites gives a segmented prefix scan followed by a short causal
transition scan.

\begin{proposition}[No bounded convolutional ledger]
No fixed finite-receptive-field operator can reproduce the zero-sign filling
rule for every line length.
\end{proposition}

\begin{proof}
Fix a receptive-field radius \(R\) and a zero secant at a site whose previous
\(R\) secants are also zero. Construct two lines with identical secants on
that complete receptive field, but place a positive witness immediately
before it in one line and a negative witness there in the other. The formal
ledger assigns opposite signs at the selected zero site. A bounded local
operator receives identical inputs and must return the same value, a
contradiction.
\end{proof}

Therefore the scan must be retained rather than approximated by a bounded
convolution. It involves linear work with one sign of persistent state per
component, while all lines of a Cartesian factor remain independent.

\subsection{Vectorized scalar-threshold admission}

CONV$^{*}$ batches the exact scalar projection
\eqref{eq:thresholdcurrent}--\eqref{eq:thresholdroot} by using a fixed sorting
network for the five breakpoints in each cell. The six resulting affine
intervals are tested in parallel, after which five signed clips return the
admitted current. Chunking consequently bounds temporary storage
independently of image size, and neither a combinatorial search nor an
iterative numerical solve is involved.

\subsection{Polyphase synthesis}

Let \(B_j^5(u)\) denote the degree-five Bernstein basis. Eliminating the
integrated controls gives
\begin{equation}
 \widehat y(i+u)=y_i+\sum_{k=0}^{4}g_k(u)c_{ik},\qquad
 g_k(u)=\sum_{j=k+1}^{5}B_j^5(u).
 \label{eq:starpolyphase}
\end{equation}
For integer refinement \(r\), only the phases \(u=p/r\),
\(p=1,\ldots,r-1\), are required. Their \(g_k\) values form a fixed
polyphase matrix. At twofold refinement,
\begin{equation}
 \widehat y(i+\tfrac12)=y_i+\frac{1}{32}
 [31,26,16,6,1]c_i.
 \label{eq:starmidpoint}
\end{equation}
Thus neither the Bezier controls nor their cumulative current sums need be
stored. Nested source nodes are copied directly for point interpolation,
whereas conservative multiresolution compiles the half-cell functional
\eqref{eq:halfcells}, block averages, and centered moment channels of
Section~IX; it does not replace them by strided anchor selection.
Accordingly, \eqref{eq:starmidpoint} is a compiled twofold phase rather than
the domain of CONV. For an arbitrary target length \(M\), the continuous
operator evaluates \eqref{eq:starpolyphase} at the coordinates
\eqref{eq:targetmap}; whenever the target uses a repeated finite phase set,
the same expression yields a fixed polyphase bank.

\subsection{Equivalence audit}

The binary64 audit uses 72 deterministic Gaussian three-channel lines over
lengths \(5,9,17,33,65,129\). The maximum difference from the formal current
is \(9.44\times10^{-16}\); the maximum synthesized-value difference is
\(8.88\times10^{-16}\). Cardinal error is zero, current-mass residual is at
most \(1.78\times10^{-15}\), and no sign-variation surplus is observed.

A second audit in Section~XI compares the local action compilation with the
formal all-source Eikonal operator. It reports the defect in
\eqref{eq:stardefect}, rather than conflating that approximation with the
binary64-equivalent FIR factor compilation.

\subsection{FIR refinement}

Collecting the preceding reductions, CONV$^{*}$ comprises compact FIR
analysis, exact segmented lineage, vectorized scalar-threshold admission,
polyphase synthesis, and the cardinal \(Q_1\) action coordinate
\eqref{eq:starcardinalcompilation}. Its one-dimensional factor is
algebraically identical to CONV; in two dimensions, the only departure is the
explicitly bounded defect \eqref{eq:stardefect}.

\subsection{Can CONV\texorpdfstring{$^{*}$}{*} Compete with EASU?}

EASU produces each target pixel in a single two-dimensional pass over a fixed
neighborhood \cite{lottes2021}. CONV$^{*}$, however, retains a barrier between
its Cartesian factors, since the second factor analyzes the nonlinear output
of the first. No throughput ordering follows from the formulas alone;
nevertheless, every nonlinear cell operation has now been reduced to the
fixed scalar-threshold and scan primitives specified above.

\subsubsection{Parallel lineage}

Sign inheritance is obtained by seeded prefix propagation, and the boundary
objective likewise has a scan form. For one transition window
\([p,q]\), let \(s^-\) and \(s^+\) be its left and right coarse signs and
define
\begin{align}
 P^-(b)&=\sum_{k=p}^{b-1}[\min(0,s^-a_k)]^2,\nonumber\\
 S^+(b)&=\sum_{k=b}^{q}[\min(0,s^+a_k)]^2.
 \label{eq:starconflictscan}
\end{align}
Using these partial sums, every candidate cost in \eqref{eq:boundary} becomes
\begin{equation}
 C(b)=P^-(b)+S^+(b),
 \label{eq:starboundarycost}
\end{equation}
so that one prefix sum and one suffix sum replace repeated evaluation of the
overlapping windows.

The dependence between successive transition boundaries is finite as well.
Indeed, if \(\kappa_\ell-\kappa_{\ell-1}\geq2\), then
\begin{equation}
 b_{\ell-1}+1\leq5\kappa_\ell-5<5\kappa_\ell-4,
 \label{eq:starwindowseparation}
\end{equation}
and the preceding boundary cannot restrict the candidate set
\eqref{eq:candidates}; dependence remains only inside a run of adjacent
coarse transitions. Writing
\(r_\ell=b_\ell-5\kappa_\ell\in\{-4,\ldots,4\}\), each transition in such a
run defines a deterministic map \(G_\ell\) from the preceding offset to the
next, including the stated first-minimum tie rule. Since function composition
is associative, a parallel prefix composition of these nine-state maps yields
the exact greedy result.

\subsubsection{Factor schedule}

The remaining work of one factor is fixed local analysis, the lineage scans,
the scalar solve \eqref{eq:thresholdmass}, and the polyphase dot product
\eqref{eq:starpolyphase}. Node jets may be formed once and shared by their two
incident cells. Fixed-scale phases are constants, and components in the
declared channel basis may be evaluated together. A tiled implementation may
write the synthesized factor through an on-chip transpose so that the next
factor again reads contiguous lines; it may not remove the intervening factor
barrier without changing the operator.

These reductions preserve the boundary closures, ordered tie rules, channel
basis, and Cartesian factor order of CONV. They reduce the remaining
nonlinear state to scans and a five-value piecewise-affine root. Whether that
arithmetic advantage offsets the intermediate factor traffic is a property
of a concrete implementation and target architecture. No performance claim
is made here; equations \eqref{eq:thresholdcurrent}--\eqref{eq:starwindowseparation}
specify the exact reduced form for such an implementation.

\end{multicols}
\twocolumn
